Table~\ref{tab:main} gives the main comparison and Fig.~\ref{fig:curves} the mean curves on two representative tasks. All statements below are about 5 seeds at $n=200$.

\begin{table*}[t]\centering\scriptsize
\caption{Main results (mean $\pm$ sd over 5 seeds; bold marks the lowest mean, not a significant difference; in Table~\ref{tab:h4}, $p_{peak}$ compares MinNorm's peak with an edge value for the ridge systems). For the ridge systems the window maximum is usually the left edge ($N/n=0.4$), not a peak. \texttt{log10\_peak\_mse}: $\log_{10}$ of the maximum test MSE in $0.4\le N/n\le4$; \texttt{bump\_rel}: hump height over best error (0 = no hump); \texttt{final\_mse}: test MSE at $N/n=25$. Task name = dataset and label-noise sd. $n$ = seeds.}\label{tab:main}
\begin{tabular}{llcccc}
\toprule
Method & Task & log10\_peak\_mse $\downarrow$ & bump\_rel $\downarrow$ & final\_mse $\downarrow$ & $n$ \\
\midrule
MinNorm & synth\_n0 & 2.151 $\pm$ 0.482 & 1634.730 $\pm$ 2421.497 & 0.162 $\pm$ 0.006 & 5 \\
FixedRidge & synth\_n0 & -0.371 $\pm$ 0.041 & 0.020 $\pm$ 0.015 & 0.370 $\pm$ 0.011 & 5 \\
GlobalRidge & synth\_n0 & \underline{-0.484 $\pm$ 0.068} & \textbf{0.006 $\pm$ 0.005} & \underline{0.155 $\pm$ 0.004} & 5 \\
\textbf{TunedRidge} & synth\_n0 & \textbf{-0.494 $\pm$ 0.061} & \underline{0.014 $\pm$ 0.006} & \textbf{0.151 $\pm$ 0.003} & 5 \\
\midrule
MinNorm & synth\_n0.25 & 2.272 $\pm$ 0.484 & 1688.381 $\pm$ 2144.831 & 0.192 $\pm$ 0.012 & 5 \\
FixedRidge & synth\_n0.25 & -0.368 $\pm$ 0.041 & \underline{0.019 $\pm$ 0.013} & 0.374 $\pm$ 0.007 & 5 \\
GlobalRidge & synth\_n0.25 & \underline{-0.467 $\pm$ 0.065} & \textbf{0.007 $\pm$ 0.008} & \underline{0.176 $\pm$ 0.010} & 5 \\
\textbf{TunedRidge} & synth\_n0.25 & \textbf{-0.470 $\pm$ 0.065} & 0.023 $\pm$ 0.015 & \textbf{0.171 $\pm$ 0.008} & 5 \\
\midrule
MinNorm & synth\_n0.5 & 2.406 $\pm$ 0.551 & 1634.050 $\pm$ 1745.121 & 0.276 $\pm$ 0.025 & 5 \\
FixedRidge & synth\_n0.5 & -0.360 $\pm$ 0.042 & \underline{0.019 $\pm$ 0.011} & 0.381 $\pm$ 0.005 & 5 \\
GlobalRidge & synth\_n0.5 & \underline{-0.425 $\pm$ 0.063} & \textbf{0.015 $\pm$ 0.013} & \underline{0.222 $\pm$ 0.021} & 5 \\
\textbf{TunedRidge} & synth\_n0.5 & \textbf{-0.428 $\pm$ 0.060} & 0.033 $\pm$ 0.024 & \textbf{0.216 $\pm$ 0.016} & 5 \\
\midrule
MinNorm & synth\_n1 & 2.761 $\pm$ 0.466 & 1488.155 $\pm$ 1268.318 & 0.608 $\pm$ 0.074 & 5 \\
FixedRidge & synth\_n1 & -0.334 $\pm$ 0.043 & \textbf{0.019 $\pm$ 0.008} & 0.405 $\pm$ 0.011 & 5 \\
GlobalRidge & synth\_n1 & \underline{-0.349 $\pm$ 0.058} & \underline{0.019 $\pm$ 0.007} & \underline{0.326 $\pm$ 0.027} & 5 \\
\textbf{TunedRidge} & synth\_n1 & \textbf{-0.350 $\pm$ 0.058} & 0.025 $\pm$ 0.014 & \textbf{0.326 $\pm$ 0.028} & 5 \\
\midrule
MinNorm & digits\_n0 & 1.797 $\pm$ 1.175 & 20675.881 $\pm$ 27747.736 & \underline{0.021 $\pm$ 0.002} & 5 \\
FixedRidge & digits\_n0 & -1.291 $\pm$ 0.014 & 0.005 $\pm$ 0.002 & 0.045 $\pm$ 0.001 & 5 \\
GlobalRidge & digits\_n0 & \underline{-1.384 $\pm$ 0.029} & \underline{0.000 $\pm$ 0.001} & 0.022 $\pm$ 0.001 & 5 \\
\textbf{TunedRidge} & digits\_n0 & \textbf{-1.404 $\pm$ 0.020} & \textbf{0.000 $\pm$ 0.000} & \textbf{0.021 $\pm$ 0.001} & 5 \\
\midrule
MinNorm & digits\_n0.15 & 2.292 $\pm$ 1.175 & 29797.484 $\pm$ 39748.856 & 0.043 $\pm$ 0.004 & 5 \\
FixedRidge & digits\_n0.15 & -1.281 $\pm$ 0.017 & \underline{0.005 $\pm$ 0.002} & 0.046 $\pm$ 0.001 & 5 \\
GlobalRidge & digits\_n0.15 & \textbf{-1.363 $\pm$ 0.025} & \textbf{0.003 $\pm$ 0.002} & \underline{0.030 $\pm$ 0.001} & 5 \\
\textbf{TunedRidge} & digits\_n0.15 & \underline{-1.359 $\pm$ 0.022} & 0.016 $\pm$ 0.027 & \textbf{0.029 $\pm$ 0.001} & 5 \\
\midrule
MinNorm & digits\_n0.3 & 2.777 $\pm$ 1.221 & 53745.601 $\pm$ 71681.344 & 0.107 $\pm$ 0.010 & 5 \\
FixedRidge & digits\_n0.3 & -1.259 $\pm$ 0.020 & \underline{0.005 $\pm$ 0.002} & 0.049 $\pm$ 0.001 & 5 \\
GlobalRidge & digits\_n0.3 & \underline{-1.287 $\pm$ 0.038} & \textbf{0.004 $\pm$ 0.002} & \textbf{0.040 $\pm$ 0.002} & 5 \\
\textbf{TunedRidge} & digits\_n0.3 & \textbf{-1.296 $\pm$ 0.026} & 0.009 $\pm$ 0.009 & \underline{0.040 $\pm$ 0.002} & 5 \\
\midrule
MinNorm & digits\_n0.6 & 3.333 $\pm$ 1.247 & 156306.900 $\pm$ 221817.677 & 0.361 $\pm$ 0.033 & 5 \\
FixedRidge & digits\_n0.6 & \textbf{-1.185 $\pm$ 0.027} & \textbf{0.005 $\pm$ 0.002} & \textbf{0.058 $\pm$ 0.004} & 5 \\
GlobalRidge & digits\_n0.6 & \underline{-1.185 $\pm$ 0.027} & \underline{0.005 $\pm$ 0.002} & \underline{0.058 $\pm$ 0.004} & 5 \\
\textbf{TunedRidge} & digits\_n0.6 & -1.185 $\pm$ 0.027 & 0.014 $\pm$ 0.021 & 0.058 $\pm$ 0.004 & 5 \\
\bottomrule
\end{tabular}

\end{table*}

\begin{figure*}[t]\centering
\includegraphics[width=0.8\textwidth]{fig_curves.pdf}
\caption{Test MSE against $N/n$ (mean of 5 seeds, log--log axes) for the four systems. The dotted line is the interpolation threshold $N=n$. MinNorm has a sharp peak there; the tuned-ridge curves decrease without a visible hump.}\label{fig:curves}
\end{figure*}

\textbf{H1 (position) -- supported.} For MinNorm the argmax of the test error is at $N/n=1$ in all \rhval{count/main/minnorm/runs} runs (Table~\ref{tab:noise}, column \texttt{pos}, all eight task medians $1.00$). Training MSE at $N=n$ is at rounding-error level (e.g.\ a mean of \rhval{main/minnorm/synth_n0/train_mse_thresh/mean} on synth\_n0; logged as \texttt{train\_mse\_thresh}), so these fits do interpolate. The grid has a resolution of $0.05$ near the threshold, so we can say only that the peak lies within one grid step of $N/n=1$.

\textbf{H2 (noise raises the peak) -- supported, with low power and a caveat.} The mean peak MSE of MinNorm increases monotonically with noise on both datasets (Table~\ref{tab:noise}), and the peak is higher at the highest than at the lowest noise level in each of the 5 paired seeds on both datasets, so the registered one-sided paired Wilcoxon test of the highest against the lowest noise level takes the smallest $p$-value attainable with 5 pairs (we do not print it: \texttt{rh compare} has no test across tasks, and the value from our own script is in \texttt{results/analysis.md}). The $\log_{10}$ peaks show the same ordering. The caveat is that the peak is already large at zero label noise (synthetic $\sigma=0$: mean peak \rhval{main/minnorm/synth_n0/peak_mse/mean} against best error \rhval{main/minnorm/synth_n0/best_mse/mean:3}). Because our teacher is not in the span of finitely many ReLU features, and the digits labels are not represented exactly by a random-feature model, approximation error may act as effective noise and produce a peak without label noise; we did not isolate this with a run. The noise dependence is much steeper on digits (Table~\ref{tab:noise}). The peak position does not move with noise.

\begin{table*}[t]\centering\small
\caption{MinNorm peak height against label noise. \texttt{pos}: median peak $N/n$; $\log_{10}$ peak: mean $\pm$ sd; mean and median of the peak test MSE over seeds; best: mean over seeds of the minimum over the grid of the test MSE.}\label{tab:noise}
\begin{tabular}{lccccc}\hline
Task & pos & $\log_{10}$ peak & mean peak & median peak & best \\ \hline
syn 0 & \rhval{main/minnorm/synth_n0/peak_pos/median:2} & \rhval{main/minnorm/synth_n0/log10_peak_mse/mean:2} $\pm$ \rhval{main/minnorm/synth_n0/log10_peak_mse/std:2} & \rhval{main/minnorm/synth_n0/peak_mse/mean} & \rhval{main/minnorm/synth_n0/peak_mse/median} & \rhval{main/minnorm/synth_n0/best_mse/mean:3} \\
syn 0.25 & \rhval{main/minnorm/synth_n0.25/peak_pos/median:2} & \rhval{main/minnorm/synth_n0.25/log10_peak_mse/mean:2} $\pm$ \rhval{main/minnorm/synth_n0.25/log10_peak_mse/std:2} & \rhval{main/minnorm/synth_n0.25/peak_mse/mean} & \rhval{main/minnorm/synth_n0.25/peak_mse/median} & \rhval{main/minnorm/synth_n0.25/best_mse/mean:3} \\
syn 0.5 & \rhval{main/minnorm/synth_n0.5/peak_pos/median:2} & \rhval{main/minnorm/synth_n0.5/log10_peak_mse/mean:2} $\pm$ \rhval{main/minnorm/synth_n0.5/log10_peak_mse/std:2} & \rhval{main/minnorm/synth_n0.5/peak_mse/mean} & \rhval{main/minnorm/synth_n0.5/peak_mse/median} & \rhval{main/minnorm/synth_n0.5/best_mse/mean:3} \\
syn 1 & \rhval{main/minnorm/synth_n1/peak_pos/median:2} & \rhval{main/minnorm/synth_n1/log10_peak_mse/mean:2} $\pm$ \rhval{main/minnorm/synth_n1/log10_peak_mse/std:2} & \rhval{main/minnorm/synth_n1/peak_mse/mean} & \rhval{main/minnorm/synth_n1/peak_mse/median} & \rhval{main/minnorm/synth_n1/best_mse/mean:3} \\
dig 0 & \rhval{main/minnorm/digits_n0/peak_pos/median:2} & \rhval{main/minnorm/digits_n0/log10_peak_mse/mean:2} $\pm$ \rhval{main/minnorm/digits_n0/log10_peak_mse/std:2} & \rhval{main/minnorm/digits_n0/peak_mse/mean} & \rhval{main/minnorm/digits_n0/peak_mse/median} & \rhval{main/minnorm/digits_n0/best_mse/mean:3} \\
dig 0.15 & \rhval{main/minnorm/digits_n0.15/peak_pos/median:2} & \rhval{main/minnorm/digits_n0.15/log10_peak_mse/mean:2} $\pm$ \rhval{main/minnorm/digits_n0.15/log10_peak_mse/std:2} & \rhval{main/minnorm/digits_n0.15/peak_mse/mean} & \rhval{main/minnorm/digits_n0.15/peak_mse/median} & \rhval{main/minnorm/digits_n0.15/best_mse/mean:3} \\
dig 0.3 & \rhval{main/minnorm/digits_n0.3/peak_pos/median:2} & \rhval{main/minnorm/digits_n0.3/log10_peak_mse/mean:2} $\pm$ \rhval{main/minnorm/digits_n0.3/log10_peak_mse/std:2} & \rhval{main/minnorm/digits_n0.3/peak_mse/mean} & \rhval{main/minnorm/digits_n0.3/peak_mse/median} & \rhval{main/minnorm/digits_n0.3/best_mse/mean:3} \\
dig 0.6 & \rhval{main/minnorm/digits_n0.6/peak_pos/median:2} & \rhval{main/minnorm/digits_n0.6/log10_peak_mse/mean:2} $\pm$ \rhval{main/minnorm/digits_n0.6/log10_peak_mse/std:2} & \rhval{main/minnorm/digits_n0.6/peak_mse/mean} & \rhval{main/minnorm/digits_n0.6/peak_mse/median} & \rhval{main/minnorm/digits_n0.6/best_mse/mean:3} \\
\hline\end{tabular}

\end{table*}

\textbf{H3 (ridge lowers the peak monotonically) -- refuted as registered.} Table~\ref{tab:sweep} and Fig.~\ref{fig:sweep} show the mean window peak against $\lambda$. It falls by orders of magnitude from $\lambda=0$ to $10^{-3}$ on every task, reaches its minimum at $10^{-3}$ or $10^{-2}$, and rises again at $\lambda=10^{-1}$ on every task, so the registered ``non-increasing in $\lambda$'' is false on all eight tasks. The rise is not a new peak: the values for $\lambda\ge1$ are nearly identical across noise levels within each dataset, consistent with a flat, under-fitted curve, and the median peak position jumps to arbitrary positions (Fig.~\ref{fig:sweep}, right); we did not run a separate check of the under-fitting explanation. With $\lambda\le10^{-6}$ the median peak position stays near $N/n=1$, for $\lambda=10^{-3}$ and $10^{-2}$ the median position is the left edge of the window ($0.4$), and for larger $\lambda$ it is at scattered interior positions (medians such as \rhval{sweep_lambda/fixedridge-sweep@lam=0.1/synth_n0/peak_pos/median:1}, \rhval{sweep_lambda/fixedridge-sweep@lam=0.1/digits_n0.15/peak_pos/median:2} and \rhval{sweep_lambda/fixedridge-sweep@lam=1/synth_n0/peak_pos/median:1}), which says only that the peak is no longer located at the threshold. Larger noise may move the best fixed $\lambda$ upward: the lowest mean peak is at $\lambda=10^{-2}$ for synth\_n1 and digits\_n0.6 and at $10^{-3}$ for five tasks, with digits\_n0.3 essentially a tie (\rhval{main/fixedridge/digits_n0.3/peak_mse/mean:4} at $10^{-2}$ against \rhval{sweep_lambda/fixedridge-sweep@lam=0.001/digits_n0.3/peak_mse/mean:4} at $10^{-3}$; Table~\ref{tab:sweep}). The statement ``the peak decreases with $\lambda$ up to $10^{-3}$ on all tasks'' was formulated after seeing the data.

\begin{table*}[t]\centering\small
\caption{Mean peak test MSE (window $0.4\le N/n\le4$) of fixed-ridge regression against $\lambda$ ($0$ = MinNorm, $10^{-2}$ = the FixedRidge of Table~\ref{tab:main}).}\label{tab:sweep}
\begin{tabular}{lcccccccc}\hline
Task & 0 & $10^{-6}$ & $10^{-4}$ & $10^{-3}$ & $10^{-2}$ & $10^{-1}$ & 1 & 10 \\ \hline
syn 0 & 269 & 1.73 & 0.377 & 0.323 & 0.427 & 0.595 & 0.631 & 0.635 \\
syn 0.25 & 325 & 2.22 & 0.415 & 0.34 & 0.43 & 0.595 & 0.631 & 0.635 \\
syn 0.5 & 441 & 4.11 & 0.561 & 0.382 & 0.438 & 0.596 & 0.631 & 0.635 \\
syn 1 & 850 & 11.9 & 1.28 & 0.546 & 0.466 & 0.599 & 0.631 & 0.635 \\
dig 0 & 458 & 0.203 & 0.0463 & 0.0395 & 0.0512 & 0.0787 & 0.0958 & 0.0995 \\
dig 0.15 & 1.32e+03 & 0.46 & 0.0576 & 0.0434 & 0.0523 & 0.0791 & 0.0959 & 0.0995 \\
dig 0.3 & 4.26e+03 & 1.19 & 0.105 & 0.0554 & 0.0551 & 0.0797 & 0.096 & 0.0995 \\
dig 0.6 & 1.63e+04 & 4.07 & 0.31 & 0.104 & 0.0654 & 0.0812 & 0.0962 & 0.0996 \\
\hline\end{tabular}

\end{table*}

\begin{figure*}[t]\centering
\includegraphics[width=0.85\textwidth]{fig_sweep.pdf}
\caption{Ridge sweep. Left: $\log_{10}$ mean peak MSE against $\log_{10}\lambda$ (leftmost point $\lambda=0$). Right: median peak position. Lines are tasks (solid: synthetic, dashed: digits; lighter = more noise).}\label{fig:sweep}
\end{figure*}

\textbf{H4 (tuned ridge removes the peak) -- partly supported.} The mean hump of TunedRidge is below $0.05$ of the best error on all eight tasks (Table~\ref{tab:h4}, column Tuned; largest task mean \rhval{main/tunedridge/synth_n0.5/bump_rel/mean:4} on synth\_n0.5), compared with task means between \rhval{main/minnorm/synth_n1/bump_rel/mean} and \rhval{main/minnorm/digits_n0.6/bump_rel/mean} for MinNorm (Table~\ref{tab:main}); the final error of TunedRidge is at or below that of MinNorm on every task (Table~\ref{tab:main}, e.g.\ synth\_n1 \rhval{main/tunedridge/synth_n1/final_mse/mean:3} against \rhval{main/minnorm/synth_n1/final_mse/mean:3}; digits\_n0 \rhval{main/tunedridge/digits_n0/final_mse/mean:3} against \rhval{main/minnorm/digits_n0/final_mse/mean:3}, equal to three decimals). Three things qualify this. (i) The hump is small but not zero: three tasks have a seed above $0.05$ (column Tuned max: synth\_n0.5, digits\_n0.15, digits\_n0.6). (ii) The second part of the registered test, that the hump is lower than MinNorm's under a paired test on \texttt{bump\_rel}, is \emph{not} met: $p$ ranges from \rhval{cmp/main/minnorm/synth_n1/bump_rel/paired_p:2} to \rhval{cmp/main/minnorm/synth_n0/bump_rel/paired_p:2} (column $p_{bump}$), because the standard deviation of MinNorm's hump is of the order of its mean. The unregistered paired $t$-test on the $\log_{10}$ peak is far below $0.05$ on all tasks (column $p_{peak}$), but we chose it after seeing the heavy tail, so it is supporting evidence and not the registered verdict. (iii) FixedRidge also has a small hump (Table~\ref{tab:main}) while being a poor predictor at low noise (synth\_n0 final MSE \rhval{main/fixedridge/synth_n0/final_mse/mean:3} against \rhval{main/minnorm/synth_n0/final_mse/mean:3} for MinNorm), so ``no peak'' alone is not evidence of good generalisation. TunedRidge, unlike FixedRidge, adapts to noise: the selected $\lambda$ grows with the noise level (Fig.~\ref{fig:lam}).

\begin{table*}[t]\centering\small
\caption{Residual hump (\texttt{bump\_rel}) of tuned-ridge variants, mean over 5 seeds. Tuned: per-width validation with 200 points (the method); Tuned max: largest seed; Global: one $\lambda$ for all widths; LOO: leave-one-out on training data; val50: 50 validation points. $p_{bump}$: paired $t$-test of Tuned against MinNorm on \texttt{bump\_rel} (registered; value from \texttt{rh compare}); $p_{peak}$: same on $\log_{10}$ peak (not registered).}\label{tab:h4}
\begin{tabular}{lccccccc}\hline
Task & Tuned & Tuned max & Global & LOO & val50 & $p_{bump}$ & $p_{peak}$ \\ \hline
syn 0 & 0.0136 & 0.0229 & 0.0065 & 0.0409 & 0.0183 & 0.21 & 0.0002 \\
syn 0.25 & 0.0226 & 0.0416 & 0.0075 & 0.0466 & 0.0501 & 0.15 & 0.0002 \\
syn 0.5 & 0.0328 & 0.0598 & 0.0147 & 0.0382 & 0.0457 & 0.10 & 0.0004 \\
syn 1 & 0.0251 & 0.0468 & 0.0193 & 0.0600 & 0.1555 & 0.06 & 0.0001 \\
dig 0 & 0.0000 & 0.0000 & 0.0005 & 0.0029 & 0.0000 & 0.17 & 0.0036 \\
dig 0.15 & 0.0155 & 0.0636 & 0.0025 & 0.0052 & 0.0271 & 0.17 & 0.0022 \\
dig 0.3 & 0.0092 & 0.0221 & 0.0043 & 0.0113 & 0.0144 & 0.17 & 0.0016 \\
dig 0.6 & 0.0139 & 0.0514 & 0.0047 & 0.0047 & 0.0318 & 0.19 & 0.0012 \\
\hline\end{tabular}

\end{table*}

\begin{figure}[t]\centering
\includegraphics[width=0.78\columnwidth]{fig_lam.pdf}
\caption{Validation-selected $\log_{10}\lambda$ against $N/n$ for TunedRidge (mean of 5 seeds), by label-noise sd.}\label{fig:lam}
\end{figure}
